\documentclass[../main]{subfiles}
\begin{document}
\chapter{Calculo de Capacitor}
\img{images/08/capacitor2}{El capacitor básico consiste en dos placas conductoras separadas por un dieléctrico no conductor que almacena energía como regiones polarizadas en el campo eléctrico entre las dos placas.}

\video{https://www.youtube.com/watch?v=h_m6qFRNITU}{Capacitores explicados}

\info{Condensador eléctrico}{
En el caso de un condensador plano, la capacidad puede expresarse por la siguiente ecuación:


A es el área efectiva de las placas;
y d es la distancia entre las placas o espesor del dieléctrico.

\begin{equation*}
   C = \epsilon \times \frac{S}{d}  
\end{equation*}

siendo:

$C:$ Capacidad del capacitor [Faradios]

$\epsilon$: Permitividad absoluta [F/m].

$d$: Distancia entre placas [m]

$S$: Superficie enfrentada entre placas [$m^2$]

}



Se denomina condensador al dispositivo formado por dos conductores cuyas cargas son iguales pero de signo opuesto.

La capacidad $C$ de un condensador se define como el cociente entre la carga $Q$ y la diferencia de potencia $V-V’$ existente entre ellos.
\info{Enería almacenada en un capacitor}{
Un condensador acumula una energía $U$ en forma de campo eléctrico. La fórmula como demostraremos más abajo es

\begin{equation*}
    U=\dfrac{1}{2} \dfrac{Q^2}{C}
\end{equation*}
}

\info{Permitividad relativa $\epsilon_r$}{

\begin{equation*}
    \epsilon = \epsilon_r \times \epsilon_0
\end{equation*}

Siendo:

$\epsilon$: Permitividad absoluta $[F/m]$.

$\epsilon_r$: Permitividad relativa (adimensional).

$\epsilon_0$: Permitividad del vacío.

\begin{equation*}
    \epsilon_0= 8,85 \times 10^{-12} \frac{F}{m}
\end{equation*}
}

\begin{table}[h]
    \centering
    \begin{tabular}{|c|r|} \hline
        Aislante &  $\varepsilon_r$ \\ \hline \hline
        Vacío & 1\\ \hline
         Aire & 1,0005\\\hline
         Nafta & 2,35\\\hline
         Aceite & 2,8\\\hline
         Vidrio & 4,7\\\hline
         Mica & 5,6\\\hline
         Glicerina & 45\\\hline
         Agua & 80,5\\ \hline
     \end{tabular}
    \caption{Permitividad relativa $\varepsilon_r$ para diferentes materiales}
\end{table}



\actividad{Resolver los siguientes problemas.}

\begin{enumerate}
    \item Calcular la permitividades absolutas para la nafta, aceite y agua.
    \item Calcular las premitividades absolutas para vidrio, mica y glicerina.
    \item Un capacitor está formado por 2 placas paralelas de $s=1m^2$ si el dieléctrico es el mica y la distancia entre placas de 0.01mm, calcular su capacidad.
    
    \item Un capacitor de 2 placas paralelas de $10cmx10cm$ si el dieléctrico es el vacío y la distancia entre placas de 0.1mm, calcular su capacidad.
    
    \item Se desea saber la capacidad de el capacitor que posee placas paralelas de $S=0,1m^2$ si el dieléctrico es el vacío $\epsilon_0$ y la distancia entre placas de 0.001mm. ¿Cómo cambia el valor de la capacidad si el dieléctrico es \emph{vidrio}?
    
    \item Calcular la energía que tiene un capacitor de $10\mu F$ si este tiene una carga de $Q=2C$

     \item Se desea saber la energía que puede almacenar un capacitor que se carga con $Q=1,8C$ y posee placas paralelas de $S=1m^2$ si el dieléctrico es el glicerina y la distancia entre placas de 0.001mm.

    \item Una placa de metal de $2 \, \text{m}^2$ se coloca a una distancia de $0.5 \, \text{mm}$ de otra placa de metal paralela de igual área. Si la permitividad del medio entre las placas es $\epsilon = 8.85 \times 10^{-12} \, \text{F/m}$, ¿Cuál es la \textbf{capacidad} del capacitor?

    \item Un capacitor plano tiene una capacidad de $100 \, \text{pF}$, y la distancia entre sus placas es de $0.1 \, \text{mm}$. Si la permitividad del material entre las placas es $\epsilon = 5 \times 10^{-11} \, \text{F/m}$, ¿Cuál es el \textbf{área} de cada placa?

    \item Un estudiante encuentra un capacitor en su laboratorio con placas de $0.01 \, \text{m}^2$ de área separadas por $0.2 \, \text{mm}$. Si la capacidad medida es de $4.425 \, \text{pF}$, ¿Cuál es la \textbf{permitividad} del material entre las placas?

    \item Dos placas de área $500 \, \text{cm}^2$ están separadas por $0.05 \, \text{mm}$ en un capacitor. Si la permitividad del material entre las placas es $\epsilon = 7 \times 10^{-12} \, \text{F/m}$, ¿Cuál es la \textbf{capacidad} del capacitor?

    \item En un experimento, se utilizan dos placas cuadradas de $0.03 \, \text{m}$ de lado, separadas por $0.01 \, \text{cm}$. Si la capacidad medida es $25 \, \text{pF}$, ¿Cuál es la \textbf{permitividad} del material entre las placas?

    \item Un capacitor tiene una capacidad de $500 \, \text{pF}$ con placas de $0.04 \, \text{m}^2$ de área. Si la distancia entre las placas es $0.2 \, \text{mm}$, ¿Cuál es la \textbf{permitividad} del material entre las placas?

    \item Dos placas de metal de $0.5 \, \text{m}^2$ de área se encuentran separadas por $1 \, \text{mm}$ en un capacitor. Si la capacidad es $22.125 \, \text{pF}$, ¿Cuál es la \textbf{permitividad} del medio entre las placas?

    \item Un capacitor con placas de $300 \, \text{cm}^2$ de área tiene una capacidad de $10 \, \text{nF}$. Si la permitividad del material entre las placas es $\epsilon = 8.85 \times 10^{-12} \, \text{F/m}$, ¿Cuál es la \textbf{distancia} entre las placas?

    \item Un experimento requiere un capacitor de $150 \, \text{pF}$. Si las placas tienen un área de $0.06 \, \text{m}^2$ y la permitividad del material es $\epsilon = 8.85 \times 10^{-12} \, \text{F/m}$, ¿Cuál debe ser la \textbf{distancia} entre las placas?

    \item En un capacitor de placas paralelas, la distancia entre las placas es de $0.01 \, \text{cm}$ y la capacidad es $500 \, \text{pF}$. Si la permitividad del material entre las placas es $\epsilon = 6 \times 10^{-12} \, \text{F/m}$, ¿Cuál es el \textbf{área} de cada placa?

    \item Un capacitor utiliza un material dieléctrico con una permitividad relativa de $\epsilon_r = 5$ y la permitividad del vacío es $\epsilon_0 = 8.85 \times 10^{-12} \, \text{F/m}$. ¿Cuál es la \textbf{permitividad} del material dieléctrico utilizado en el capacitor?

    \item En un experimento, se mide la permitividad de un material dieléctrico como $\epsilon = 2.655 \times 10^{-11} \, \text{F/m}$. Si la permitividad del vacío es $\epsilon_0 = 8.85 \times 10^{-12} \, \text{F/m}$, ¿Cuál es la \textbf{permitividad relativa} $\epsilon_r$ del material?

    \item Un capacitor tiene un material dieléctrico con una permitividad de $\epsilon = 1.77 \times 10^{-11} \, \text{F/m}$ y una permitividad relativa de $\epsilon_r = 2$. ¿Cuál es la \textbf{permitividad del vacío} $\epsilon_0$?

    \item En un dispositivo electrónico, se necesita usar un material con una permitividad relativa $\epsilon_r = 10$. Si la permitividad del vacío es $\epsilon_0 = 8.85 \times 10^{-12} \, \text{F/m}$, ¿Cuál será la \textbf{permitividad} $\epsilon$ del material necesario?

    \item Un capacitor está diseñado con un material dieléctrico cuya permitividad es $\epsilon = 7.08 \times 10^{-11} \, \text{F/m}$. Si la permitividad del vacío es $\epsilon_0 = 8.85 \times 10^{-12} \, \text{F/m}$, ¿Cuál es la \textbf{permitividad relativa} $\epsilon_r$ del material utilizado en el capacitor?

\end{enumerate}

\end{document}